\documentclass[12pt]{article}
\usepackage[utf8]{inputenc}
\usepackage[T1]{fontenc}
\usepackage{hyperref}
\title{Yaeyama Language Resources for K-12 Curriculum Development}
\author{Compiled by Claude Sonnet 5.5 (claude-sonnet-5-5) for Andr\'as Kornai}
\date{}
\begin{document}
\maketitle

\section{Introduction}
Yaeyama (ISO 639-3 rys, Glottocode yaey1239) is a Japonic language of the southern Ryukyu (Yaeyama) Islands, including Ishigaki. It is endangered and is not known to be taught in schools. The search found very few Yaeyama-specific materials. This dossier lists only what appeared in search results and does not claim more.

\section{Primary Educational Resources}
No Yaeyama school textbook, primer or graded reader turned up in the searches. The Meeramuni project below is the closest thing to learner material, and it covers only the Miyara variety of Ishigaki.
\begin{itemize}
\item Meeramuni: a learning project for the Miyara variety of Yaeyaman (Ishigaki). The Endangered Languages Project describes lessons in vocabulary and grammar, a dictionary, conjugation tables and a tool for sound correspondences with Japanese. \url{https://www.endangeredlanguages.com/node/110363}
\end{itemize}

\section{Linguistic and Reference Materials}
\begin{itemize}
\item Hiroshima University graduate school bulletin article (2004) on early postwar Japanese-language reading textbooks used in Yaeyama. It concerns standard Japanese literacy and not the Yaeyama language itself, but it is useful for the history of schooling in the islands. \url{https://ir.lib.hiroshima-u.ac.jp/00018204}
\item Endangered Languages Project sample entry for Yaeyama (language page 5353). It notes that most speakers in the southern Yaeyama Islands are over 70. \url{https://endangeredlanguages.com/lang/5353/samples/11760}
\end{itemize}

\section{Religious and Community Resources}
No religious or community-organisation resources for Yaeyama were found in the searches.

\section{Digital and Audio Resources}
The Meeramuni site is the only digital resource found (see above). No audio or video collections were located.

\section{Curriculum Development Context}
Speakers are mostly elderly, and the language is not known to be taught in schools. The available material is thin and varies by island and village, so a curriculum would need new documentation work: recording elders, choosing a writing system, and building a reader and word list for each variety. The Okinawa Prefecture work on a common kana orthography for the five Ryukyuan areas, which includes Yaeyama, is relevant. It was seen in a search summary of the prefectural document at \url{https://www.pref.okinawa.lg.jp/_res/projects/default_project/_page_/001/009/625/01_gaiyou.pdf}.
\end{document}